\section{Statistical formulas and post-hoc protocols}
\label{sec:appendix_methodology}

Formal definitions the main body cites but does not display.

\subsection{Judge--human agreement: Cohen's $\kappa$ and Krippendorff's $\alpha$}
\label{sec:appendix_kappa_alpha}

For each binarised rating (judge threshold $J \geq \rho$ versus $J < \rho$),
Cohen's $\kappa$ is

\begin{equation}
  \kappa
  \;=\;
  \frac{p_o - p_e}{1 - p_e},
  \label{eq:kappa}
\end{equation}

where $p_o$ is the observed agreement proportion between the judge SLM
and a human annotator and $p_e$ is the chance-expected proportion under
independence. Values follow Landis and Koch's bands: $\kappa > 0.6$
substantial, $\kappa > 0.8$ almost perfect. Effect-size conventions
throughout the paper follow \citet{Cohen2013StatPower}. For the
multi-annotator setting (five human raters $+$ one judge), we
additionally compute Krippendorff's $\alpha$~\citep{Marzi2024KrippendorffAlpha}
for ordinal ratings $r_{ij}$ from rater $i$ on item $j$:

\begin{equation}
  \alpha
  \;=\;
  1 \;-\; \frac{D_o}{D_e},
  \quad
  D_o \;=\; \tfrac{1}{n}\sum_{j} \tfrac{1}{\binom{m_j}{2}}
            \sum_{i_1 < i_2} \delta(r_{i_1 j}, r_{i_2 j})^2,
  \label{eq:alpha}
\end{equation}

where $\delta$ is the ordinal-difference metric and $D_e$ is the
expected disagreement under random pairing. Values $\alpha > 0.667$
are considered acceptable for inferential use. Both $\kappa$ and
$\alpha$ are reported in Section~\ref{sec:results_human_eval}; the
headline numbers on the 60-pair five-annotator sample are
$\alpha_{\textsf{ord}} = 0.193$ (5-way, all annotators) and
$\kappa_{\textsf{quad}} = 0.074$ (DeepSeek-R1 judge vs.\ 5-way
majority human rating).

\subsection{Post-processing categorization of every closure row}
\label{sec:appendix_postproc}
\label{sec:experiment_setup_postproc}

Every closure row in the results TSV is post-hoc labeled with the
Algorithm~\ref{alg:validity} decision reason and the
Algorithm~\ref{alg:test_filter} filter reason. The distribution of
decision reasons per path (Table~\ref{tab:disagreement_by_path}) is
the primary judge--metric disagreement diagnostic; the distribution of
filter reasons per (cell, path) triple is the primary diagnostic for
the phi-4-in-test-stage attribution (Section~\ref{sec:discussion_rq3}).

\subsection{Capability-preservation protocol}
\label{sec:appendix_capability_protocol}
\label{sec:experiment_setup_capability}

For each path, we identify the strongest mono baseline (M4 for
Paths~1 and~3; M3 for Path~2) and restrict analysis to samples on
which that baseline achieves valid closure. For every other cell, the
\emph{capability-preservation rate} is the fraction of those same
samples on which the cell also achieves valid closure. Values near~$1$
indicate that the cell does not sacrifice mono strengths; substantially
lower values expose cells where specialization costs capability on
samples the strongest mono handles cleanly. Results are reported in
Table~\ref{tab:capability_preservation} and discussed in
Section~\ref{sec:discussion_capability}.
